\ticket{6}{Логическая модель представления знаний}

\begin{examplan}
    \item Дать определение логической модели представления знаний.
    \item Описать применение логики предикатов первого порядка и этапы формализации.
    \item Назвать основные трудности логического представления знаний.
    \item Объяснить дескриптивные логики: концепты, роли, TBox/ABox и связь с логикой предикатов.
\end{examplan}

\subsection*{Короткая версия для письменного ответа}

\textbf{Логическая модель представления знаний} описывает знания в виде формул некоторой формальной логики. Знания задаются как утверждения об объектах, их свойствах и отношениях, а новые знания получаются с помощью правил вывода.

Компоненты логической модели: формальный язык, семантика и механизм вывода.

\subsection*{Логика предикатов первого порядка}

Логика предикатов первого порядка позволяет описывать:
\begin{itemize}
    \item объекты через константы;
    \item переменные, пробегающие объекты предметной области;
    \item свойства и отношения через предикаты;
    \item функциональные зависимости через функциональные символы;
    \item общие и частные утверждения через кванторы $\forall$ и $\exists$.
\end{itemize}

Примеры:
\[
\forall x(Student(x)\to Smart(x)),
\]
\[
\forall x\forall y\forall z((Parent(x,y)\land Parent(y,z))\to Ancestor(x,z)).
\]

\subsection*{Этапы представления знаний}

\begin{enumerate}
    \item Анализ предметной области: выделение объектов, свойств, отношений и типовых утверждений.
    \item Выбор словаря: константы, предикаты, функции.
    \item Формализация знаний: перевод естественно-языковых утверждений в логические формулы.
    \item Построение базы знаний из фактов и общих правил.
    \item Формулировка запросов и запуск вывода.
\end{enumerate}

\subsection*{Основные сложности}

\begin{itemize}
    \item \textbf{Выразительность против вычислимости}: полная логика предикатов выразительна, но автоматический вывод в ней сложен и в общем случае неразрешим.
    \item \textbf{Формализация естественного языка}: неформальные знания часто контекстны и неоднозначны.
    \item \textbf{Монотонность}: в классической логике добавление фактов не отменяет старые выводы, что плохо подходит для рассуждений по умолчанию.
    \item \textbf{Неполнота и неопределенность}: классическая логика плохо выражает вероятностные и нечеткие знания без расширений.
    \item \textbf{Масштабируемость}: вывод на больших базах знаний может быть дорогим.
\end{itemize}

\subsection*{Дескриптивные логики}

\textbf{Дескриптивные логики} --- семейство разрешимых фрагментов логики предикатов, ориентированных на описание понятий предметной области и отношений между ними. Они лежат в основе OWL и многих онтологических систем.

Основные элементы:
\begin{itemize}
    \item \textbf{индивиды} --- конкретные объекты;
    \item \textbf{концепты} --- классы индивидов, аналоги унарных предикатов;
    \item \textbf{роли} --- бинарные отношения между индивидами.
\end{itemize}

Конструкторы концептов включают пересечение $C\sqcap D$, объединение $C\sqcup D$, отрицание $\neg C$, ограничения существования $\exists R.C$, всеобщности $\forall R.C$ и числовые ограничения.

\subsection*{TBox и ABox}

\begin{description}
    \item[TBox] терминологическая часть онтологии: определения и включения концептов, например $Man \equiv Person\sqcap Male$ или $Student\sqsubseteq Person$.
    \item[ABox] утверждения об индивидах, например $Person(John)$ или $hasChild(John,Mary)$.
\end{description}

Типовые задачи вывода: проверка выполнимости концепта, включения концептов, принадлежности индивида концепту и извлечение всех экземпляров концепта.

\subsection*{Связь с логикой предикатов}

Дескриптивные логики переводятся в логику предикатов:
\begin{itemize}
    \item концепт $C$ соответствует унарному предикату $C(x)$;
    \item роль $R$ соответствует бинарному предикату $R(x,y)$;
    \item $\exists R.C$ соответствует $\exists y(R(x,y)\land C(y))$;
    \item $C\sqsubseteq D$ соответствует $\forall x(C(x)\to D(x))$.
\end{itemize}

Главное преимущество дескриптивных логик --- ограничение выразительности ради разрешимого и практически применимого вывода.

\begin{important}
    \item Логические модели дают строгую семантику и формальный вывод.
    \item Логика предикатов выразительна, но сложна для автоматического вывода.
    \item Дескриптивные логики описывают концепты, роли и индивиды.
    \item TBox хранит терминологию, ABox --- факты об экземплярах.
    \item ДЛ являются фрагментами логики предикатов с контролируемой выразительностью.
\end{important}

